\chapter{Открытые вопросы}

Решённое отделено от нерешённого честно. Ниже вопросы, по которым механика ясна, а политика или детали нет, и вопросы, по которым нет и механики. Для каждого текущее состояние мысли, чтобы новый участник не начинал с нуля. С версии 0.4 там, где вопрос перешёл в спецификацию модуля, в конце раздела сказано, где он живёт сейчас.

\section{Пересечение областей}

Несколько оркестраторов с дизъюнктными областями работают. Задача, затрагивающая две области, принадлежит аренде, покрывающей объединение, либо создаётся кор-нодой напрямую. Не решено: кто и по какому правилу решает, что область нужно объединить, и как часто это будет происходить на реальном монорепозитории. Ожидание: редко, если декомпозиция идёт по модулям с явными интерфейсами; проверяется на данных.

\term{Где вопрос сейчас.} Владение областями задают аренды областей waggle (\texttt{spec/waggle.md} v5.2, §12): ядро проверяет аренду при запуске работы, при интеграции и при записи борда, а область переходит к другому оркестратору через \texttt{coord.handoff-offer} с большей эпохой аренды. Для нескольких операторов на одном борде решён один принимающий писатель с fencing по эпохе (трек CM2); как он выбирается, по владению процессом или выбором участников, открыто (вики, \texttt{open-questions.md}).

\section{Стимулы ревью}

Вестинг для одобрившего ревьювера толкает к отклонению на всякий случай, и пропускная способность падает. Правило неоплаты проигранных в споре отклонений давит в обратную сторону. Это классическая задача peer prediction, у неё нет закрытого решения, и в протокол оно не заложено. Измеряется долей пересмотренных отклонений; правила оплаты раундов настраиваются по данным.

\term{Где вопрос сейчас.} Измерять пропуски ревьювера предлагает калибровка посевом дефектов с охранным условием \texttt{review.recall} (curator-playbook, \texttt{spec/process-configuration.md} v4.8, §6.16). Лишние отклонения она не измеряет, и вопрос о стимулах остаётся открытым.

\section{Семантически вредная задача внутри своей области}

Ограничена областью, бюджетом, независимым ревью, вестингом и остановкой, но не устранена. Возможные усиления: детерминированное правило для дифов деструктивной формы (удаление больше заданной доли строк, удаление файлов вне объявленной области) с требованием дополнительного одобрения кор-ноды; проверка семантического соответствия контракта задачи родительской цели вторым ревью для обычного класса. Версионирование контракта закрывает тихое ослабление критериев, но не вредную задачу с честными критериями. Оба усиления удорожают цикл, и вопрос о пороге открыт.

\section{Перенос одобрения между кандидатами}

Базовая гарантия требует повторного ревью после каждой пересборки кандидата. Расширение вне гарантии допускает перенос при условии, что дельта базы не пересекается ни с write scope задачи, ни с множеством путей, прочитанных продьюсером по receipts. Не решено: достаточно ли множества прочитанных путей как приближения контекста ревью, и как часто условие будет выполняться на реальных очередях. Проверяется на данных пилота: доля пересборок, при которых условие выполнено, и доля дефектов, найденных повторным ревью там, где оно было бы пропущено.

\section{Механическая проверка}

Исполняемая модель фрагмента интеграции проверена перебором в заданных границах. Не сделано: та же модель в языке спецификаций с механической проверкой, распространение на аренду, остановку и внешние эффекты, и доказательство свойств для всех допустимых исполнений, а не для ограниченного пространства сценариев.

\term{Где вопрос сейчас.} Исследование формальных методов записано в роадмап экосистемы (вики, \texttt{roadmap/ecosystem-roadmap.ru.md}, §5: проверка моделей для критичных протоколов, например аренд waggle), а блок проверяемых свойств процесса с экспортом в TLA+ или Quint решено делать после PC1 (curator-playbook v4.8, §17.3, п.~17). Модель фрагмента интеграции в эти планы пока не включена.

\section{Аттестация среды}

Среда доверена в отношении целостности. На Linux доступны аппаратно аттестованные анклавы и microVM, на macOS нет ни того, ни другого для произвольной нагрузки. Верификатор на второй машине переводит тесты из подписи под ответом той же машины в независимое исполнение, но не закрывает подмену дифа, показанного ревьюверу. Возможный ход: ревьювер получает диф не из санбокса среды, а по хешу коммита из собственной копии объектов git, что сводит доверие к целостности хешей. Не спроектировано.

\section{Экономика}

Учёт вклада описан правилами, но ставки (цена единицы сложности, множители для ёмкости и для человеческого времени, доля ретроактивного пула) не назначены, а форма конвертации (репутация, деньги, доля) зависит от проекта. Для доверенной группы задача решаема как контракт плюс время плюс ограниченная дискреция. Для открытого анонимного роя пропорциональные доли не собираются в принципе, и там остаются репутация и баунти. Экономика вынесена в отдельную работу.

\section{Пул сред}

Одна авторитетная среда и реплики по снимкам описаны, но протокол синхронизации снимков между репликами и правило, какая реплика обслуживает какой санбокс, не специфицированы. Ожидание: реплики получают снимок дерева на момент назначения по хешу коммита из общего хранилища объектов, и любая реплика может обслуживать любой санбокс; кандидат строится и голова продвигается только на авторитетной.

\section{Федерация и реестр}

Манифест проекта как источник плана контрибьютора описан. Реестр манифестов, взвешивание чужих аттестаций репутацией выдавшего проекта и правила деления ёмкости между проектами описаны на уровне принципа. Всё это нужно, когда проектов больше одного, и не раньше.

\term{Где вопрос сейчас.} Общение с другими организациями есть этап waggle CM4: CA партнёров, инспекторы сообщений по умолчанию, федерация XMPP, шлюз A2A, квоты (\texttt{spec/waggle.md} v5.2, §17). Допуск внешних участников есть стадия L2 в curator-trust (\texttt{spec/trust.md} v0.4, §13): репутация удешевляет проверку, но доверия не даёт. Реестр манифестов и деление ёмкости между проектами ни одна спецификация пока не описывает.

\section{Идентичность: вес псевдонимов и хрупкость провайдеров}

Механика доказательств идентичности ясна, политика нет. Не решено: какой вес и какую глубину ревью давать псевдонимным нодам, чтобы не отрезать предложение ёмкости и не открыть дешёвые множественные личности; как часто переподтверждать доказательства, когда провайдер закрыл доступ к API или удалил публикацию; считать ли ключи аккаунта на хостинге кода доказательством того же класса, что домен под DNSSEC. Ожидание: веса настраиваются по данным первых проектов, а хрупкость провайдеров гасится требованием двух доказательств разных классов для веса выше псевдонимного.

\term{Где вопрос сейчас.} Классы ключей и уровни их гарантии (обычный, привязанный к устройству, аттестованный) задаёт waggle (\texttt{spec/waggle.md} v5.2, §4.5). Кто считается человеком, какие ключи требуют его присутствия и как ключи одного человека ручаются друг за друга, задаёт curator-trust (\texttt{spec/trust.md} v0.4, §3). Доказательства идентичности через резолверы и их веса в эти спецификации пока не перенесены.

\section{Что не решено намеренно}

Полная децентрализация без единой среды: техническая возможна через аттестованное или повторное исполнение, но не нужна, пока доверие к среде решается социально, выбором проекта. Блокчейн как расчётный слой: имеет смысл только когда проекты начнут платить контрибьюторам деньгами или потребуется стейк как барьер вместо цены ёмкости. Оба вопроса отложены с формулировкой, а не забыты.
